% Owner: Prisca. Sources: stubs/README.md, stubs/js/*/*.meta.json, docs/methodology.md, verifier/run_stub.js.
\section{Dataset and Verifier}
\label{sec:dataset}

\subsection{Stubs}
The test set has 102 hand-written JavaScript request handlers (``stubs''): 17 vulnerable/safe pairs for each of
three classes, cross-site scripting (XSS), command injection (CMDi) and path traversal (Table~\ref{tab:dataset}).
Each stub is an Express-style CommonJS module of about ten lines (median 10) that reads one input from the query,
body, route parameters or headers and passes it to a sink. The two programs in a pair are identical except for the
defense: an escaping helper, an allowlist, a regular expression, \texttt{execFile} with an argument array, a
\texttt{path.basename} or prefix check, and so on (12--14 distinct defenses per class). Comments are kept
neutral and never mention security, so the text does not leak the label. Every stub's label is confirmed by the verifier
(below). We also wrote a review checklist for each pair: the vulnerable version is exploitable, the safe version
is safe beyond our test payloads, the two differ only by the defense, and nothing hints at the label. One author
reviewed 26 of the 51 pairs by hand, after a tool-assisted first pass: all 26 are correct for their class, but
five command-injection safe stubs are debatable (for example, \texttt{execFile} with an argument array stops shell
injection but may still allow argument injection). We kept them unchanged, since all runs use them, and discuss them
in Section~\ref{sec:limitations}. \todo{status of the other 25 pairs (Eileen's review)}

\begin{table}[t]
\centering
\caption{The 102 stubs. Every class has 17 vulnerable/safe pairs.}
\label{tab:dataset}
\footnotesize
\begin{tabular}{@{}p{0.75cm}p{3.6cm}p{3.6cm}@{}}
\toprule
Class & Sinks (stubs) & Example defenses in the safe twin \\
\midrule
XSS  & \texttt{res.send} (32), \texttt{res.write}/\texttt{end} (2) & HTML escaping, allowlist, numeric coercion \\
CMDi & \texttt{exec} (22), \texttt{execSync} (10), \texttt{spawn} with shell (2) & regex allowlist, \texttt{execFile}, quoting \\
Path & \texttt{readFileSync} (18), \texttt{readFile} (10), streams/\texttt{sendFile} (6) & \texttt{basename}, resolve + prefix check \\
\bottomrule
\end{tabular}
\end{table}

\subsection{Execution-based verifier}
The verifier (\texttt{verifier/run\_stub.js}) loads a program in a Node.js \texttt{vm} context in which
\texttt{child\_process}, \texttt{fs} and the response object are replaced by recording mocks. It calls the handler
with an attack payload placed in all four input locations and checks the arguments that reach the sink:
\begin{itemize}
\item \textbf{XSS:} the response contains the raw \texttt{<script>} payload.
\item \textbf{CMDi:} with any of three payloads (plain, single-quote breakout, double-quote breakout), the injected
      \texttt{echo} follows an \emph{unquoted} shell separator. The check follows shell quoting, and the mocks reject
      argument types that real Node.js rejects.
\item \textbf{Path traversal:} the resolved path escapes the handler's base directory.
\end{itemize}
A program that crashes is reported as unverifiable, never as safe. All 102 stubs match their labels.

The verifier is what makes the results trustworthy, and it needed care. Reading the first QA edits by hand
revealed two verifier false negatives: a logging edit that wrapped the input in unescaped quotes, and a refactor
that called \texttt{spawn} with an invalid argument (which throws in Node.js). Both had been counted as fixes. After
we added the quote-breakout payloads and the argument checks, we re-verified all stored edits without new model
calls; none of the 306 outcomes of the main QA run changed. Every generated variant in Parts 2--4 is also run
through the verifier, and all keep their label.
